% Software Requirements Specification: LaTeX template
% Compile with pdfLaTeX (Overleaf default). Put figures in ./figures as PDF or PNG.

\documentclass[11pt,a4paper]{report}

% ---------- Packages ----------
\usepackage[utf8]{inputenc}
\usepackage[T1]{fontenc}
\usepackage{lmodern}
\usepackage[margin=2.5cm]{geometry}
\usepackage{graphicx}
\usepackage{float}
\usepackage{booktabs}
\usepackage{tabularx}
\usepackage{longtable}
\usepackage{array}
\usepackage{enumitem}
\usepackage[table]{xcolor}
\usepackage{tcolorbox}
\usepackage{titlesec}
\usepackage{fancyhdr}
\usepackage{caption}
\usepackage[hidelinks]{hyperref}

\graphicspath{{figures/}}
\setlength{\parindent}{0pt}
\setlength{\parskip}{6pt}

% ---------- Document details ----------
\newcommand{\productname}{Product Name}
\newcommand{\projectname}{Project Name}
\newcommand{\coursename}{Course Code: Course Title}
\newcommand{\institution}{Institution}
\newcommand{\docversion}{0.1}
\newcommand{\docdate}{\today}

% ---------- Headings, header and footer ----------
\titleformat{\chapter}{\normalfont\LARGE\bfseries}{\thechapter.}{0.6em}{}
\titlespacing*{\chapter}{0pt}{0pt}{18pt}

\pagestyle{fancy}
\fancyhf{}
\fancyhead[L]{\small \productname\ -- SRS}
\fancyhead[R]{\small Version \docversion}
\fancyfoot[C]{\thepage}
\fancypagestyle{plain}{\fancyhf{}\fancyfoot[C]{\thepage}\renewcommand{\headrulewidth}{0pt}}

% ---------- Reusable blocks ----------
% User story with confirmation criteria:
% \begin{userstory}{ID}{Title}{Priority} ... \end{userstory}
\newtcolorbox{userstory}[3]{
  colback=white, colframe=black!60, boxrule=0.5pt, arc=2pt,
  fonttitle=\bfseries, title={#1 \textperiodcentered\ #2 \hfill #3}
}

% Placeholder text, shown in red so nothing unfinished is missed
\newcommand{\todo}[1]{\textcolor{red}{[#1]}}

\begin{document}

% ---------- Title page ----------
\begin{titlepage}
  \centering
  \vspace*{3cm}
  {\Large \institution\par}
  \vspace{0.5cm}
  {\large \coursename\par}
  \vspace{2.5cm}
  {\Huge\bfseries Software Requirements Specification\par}
  \vspace{0.8cm}
  {\LARGE \productname\par}
  \vspace{0.3cm}
  {\large \projectname\par}
  \vspace{2.5cm}
  {\large Version \docversion\par}
  {\large \docdate\par}
  \vfill
  \begin{tabular}{ll}
    \toprule
    \textbf{Name} & \textbf{ID} \\
    \midrule
    \todo{Member} & \todo{ID} \\
    \todo{Member} & \todo{ID} \\
    \todo{Member} & \todo{ID} \\
    \todo{Member} & \todo{ID} \\
    \bottomrule
  \end{tabular}
\end{titlepage}

% ---------- Revision history ----------
\chapter*{Revision History}
\begin{tabularx}{\textwidth}{l l X l}
  \toprule
  \textbf{Version} & \textbf{Date} & \textbf{Changes} & \textbf{By} \\
  \midrule
  \todo{0.1} & \todo{Date} & \todo{Change} & \todo{Name} \\
  \bottomrule
\end{tabularx}

\tableofcontents
\listoffigures
\listoftables

% ---------- 1 Introduction ----------
\chapter{Introduction}
\section{Purpose}
\section{Intended Audience}
\section{Intended Use}
\section{Product Scope}
\section{Risk Definitions}

% ---------- 2 Overall Description ----------
\chapter{Overall Description}
\section{User Classes and Characteristics}
\section{User Needs}
\section{Operating Environment}
\section{Constraints}
\section{Assumptions}

% ---------- 3 Requirements ----------
\chapter{Requirements}
\section{Functional Requirements}
% Example:
% \begin{userstory}{XX-01}{Story title}{Must}
%   As a \todo{role}, I want \todo{goal}, so that \todo{benefit}.
%   \textbf{Confirmation}
%   \begin{itemize}
%     \item Given \todo{context}, when \todo{action}, then \todo{result}.
%   \end{itemize}
% \end{userstory}

\section{Non-Functional Requirements}
% Example:
% \begin{longtable}{p{2cm} p{2.4cm} p{7.6cm} p{2.2cm}}
%   \toprule
%   \textbf{ID} & \textbf{Category} & \textbf{Requirement} & \textbf{Story} \\
%   \midrule
%   \endhead
%   \todo{NFR-01} & \todo{Category} & \todo{Requirement} & \todo{IDs} \\
%   \bottomrule
% \end{longtable}

% ---------- Appendices ----------
\appendix
\chapter{Glossary}
\chapter{Analysis Models}
\section{Use Case Diagrams}
\section{Use Case Narratives}
\section{State Diagrams}
\section{Data Flow Diagrams}
% Example figure:
% \begin{figure}[H]
%   \centering
%   \includegraphics[width=\textwidth]{diagram-name}
%   \caption{Caption}
%   \label{fig:diagram-name}
% \end{figure}

\end{document}
